\chapter{Quick Reference}
\label{app:quick}

\section{Document skeleton}

\begin{latexcode}{Minimal document}
\documentclass{article}
% packages and commands
\begin{document}
Hello, world!
\end{document}
\end{latexcode}

\section{Structure and text}

\begin{center}
\begin{tabularx}{\linewidth}{@{}lX@{}}
\toprule
Source & Result or purpose \\
\midrule
\cmd{title}\required{text} & Define the document title. \\
\cmd{author}\required{name} & Define the author. \\
\cmd{maketitle} & Print title information. \\
\cmd{section}\required{text} & Numbered section. \\
\cmd{subsection}\required{text} & Numbered subsection. \\
\cmd{section*}\required{text} & Unnumbered section. \\
\cmd{emph}\required{text} & Contextual emphasis. \\
\cmd{textbf}\required{text} & Bold label or short emphasis. \\
\cmd{footnote}\required{text} & Footnote. \\
\code{\% \& \_ \$ \# \{ \}} & Common literal reserved characters. \\
\bottomrule
\end{tabularx}
\end{center}

\section{Lists}

\begin{latexcode}{List patterns}
\begin{itemize}
  \item Unordered item
\end{itemize}

\begin{enumerate}
  \item Ordered item
\end{enumerate}

\begin{description}
  \item[Term] Explanation
\end{description}
\end{latexcode}

\section{Mathematics}

\begin{center}
\begin{tabularx}{\linewidth}{@{}lX@{}}
\toprule
Source & Purpose \\
\midrule
\code{\textbackslash(...\textbackslash)} & Inline mathematics. \\
\code{\textbackslash[...\textbackslash]} & Unnumbered display. \\
\code{equation} & Numbered single equation. \\
\code{align} & Aligned numbered equations. \\
\cmd{frac}\required{a}\required{b} & Fraction. \\
\cmd{sqrt}\required{x} & Square root. \\
\code{x\_i\^{}2} & Subscript and superscript. \\
\cmd{text}\required{words} & Ordinary words inside mathematics. \\
\cmd{label}\required{eq:key} & Label a numbered equation. \\
\cmd{cref}\required{eq:key} & Type-aware reference with \pkg{cleveref}. \\
\bottomrule
\end{tabularx}
\end{center}

\section{Quantities}

\begin{latexcode}{siunitx patterns}
\num{12345.6}
\si{\milli\metre}
\qty{12.5}{\milli\gram}
\qtyrange{18}{22}{\degreeCelsius}
\end{latexcode}

\section{Tables}

\begin{latexcode}{booktabs pattern}
\begin{tabular}{lrr}
\toprule
Group & Mean & SD \\
\midrule
A & 12.4 & 1.8 \\
B & 16.7 & 2.1 \\
\bottomrule
\end{tabular}
\end{latexcode}

\section{Figures}

\begin{latexcode}{Figure pattern}
\begin{figure}[htbp]
  \centering
  \includegraphics[width=.8\linewidth]{figures/result.pdf}
  \caption{An informative caption.}
  \label{fig:result}
\end{figure}
\end{latexcode}

\section{References and citations}

\begin{center}
\begin{tabularx}{\linewidth}{@{}lX@{}}
\toprule
Source & Purpose \\
\midrule
\cmd{label}\required{key} & Attach a stable name. \\
\cmd{ref}\required{key} & Print the number. \\
\cmd{cref}\required{key} & Print the type and number. \\
\cmd{textcite}\required{key} & Author in the sentence. \\
\cmd{parencite}\required{key} & Parenthetical citation. \\
\cmd{printbibliography} & Print cited records. \\
\bottomrule
\end{tabularx}
\end{center}

\section{Build commands}

\begin{terminalcode}{Build reference}
latexmk -pdf main.tex
latexmk -C main.tex
pdfinfo main.pdf
pdffonts main.pdf
\end{terminalcode}
